\documentclass[aspectratio=169,notheorems]{corvinusmetropolis}
\usepackage{qrcode,adjustbox} % iff qrcode needed, otherwise comment this line!
\usepackage{amsmath, unicode-math}
\newtheorem{theorem}{Theorem}
\newtheorem{proposition}[theorem]{Proposition}
\newtheorem{lemma}[theorem]{Lemma}
\theoremstyle{definition}
\newtheorem{definition}[theorem]{Definition}
\newtheorem{corollary}[theorem]{Corollary}
\newtheorem{example}[theorem]{Example}
\email{magyarkuti@uni-corvinus.hu}

\author[Magyarkuti]{Gyula~Magyarkuti}
\institute[Corvinus University]
  {
  Department of Mathematics\\
  Corvinus University of Budapest \\
  S208/B

  \begin{minipage}[b]{0.8\linewidth}
  email: 
  \texttt{magyarkuti@uni-corvinus.hu}\\
  web:
  \url{https://magyarkuti.github.io/linear-algebra}
  \end{minipage}
  \hfill{\adjustbox{valign=B}{\qrcode[hyperlink]{https://magyarkuti.github.io/linear-algebra}}}%% iff qrcode needed, otherwise comment this line!
  }
\titlegraphic{\hfill{\includegraphics[width=.2\textwidth]{corvinus_Department_of_Mathematics_logo_cmyk.eps}}}
\subject{Linear Algebra}
\newcommand{\R}{\mathbb{R}}
\newcommand{\spn}{\operatorname{span}}

\renewcommand{\mathbf}{\mathbb}
\newcommand{\lin}{\operatorname{lin}}
\newcommand{\gen}{\operatorname{gen}}
\newcommand{\im}{\operatorname{Im}}
\newcommand{\id}{\operatorname{id}}
\newcommand{\sgn}{\operatorname{sgn}}
\newcommand{\tr}{\operatorname{tr}}
\newcommand{\codim}{\operatorname{codim}}
%newcommand{\deg}{\operatorname{deg}}
\newcommand{\ip}[2]{\langle#1,#2\rangle}
\newcommand{\uj}{\text{új}}
\newcommand{\rgi}{\text{régi}}
\newcommand{\Star}[1]{#1\ensuremath{^*}\kern-\scriptspace}
\newcommand{\CStar}{\Star{\ensuremath{\mathrm{C}}}}

\renewcommand{\mathbf}{\symbf}
\renewcommand{\mathcal}{\symcal}
\endinput

\title{Matrix factorization}
\subtitle{The rank of a matrix}
\date{October 7, 2024}

\begin{document}

\frame{\titlepage}
%\section*{Outline:}
\frame{\frametitle{Outline}
\tableofcontents}


\section{System of linear equations}
\frame{\frametitle{System of linear equations}
\begin{definition}\label{lsystem}
The general form of a system of linear equations is
\begin{eqnarray*}
\alpha_{11}x_1+\alpha_{12}x_2+\cdots +\alpha_{1n}x_n&=&\beta_1\\
\alpha_{21}x_1+\alpha_{22}x_2+\cdots +\alpha_{2n}x_n&=&\beta_2\\
\vdots & & \\
\alpha_{m1}x_1+\alpha_{m2}x_2+\cdots +\alpha_{mn}x_n&=&\beta_m
\end{eqnarray*}
where $\alpha_{ij}\ i=1,2,\ldots m,\ j=1,2,\ldots ,n$ and $\beta_i\ i=1,2\ldots ,m$ are given scalars and $x_j\ j=1,2,\ldots ,n$ are the unknowns.
\end{definition}
The system of linear equations is said to be \alert{homogeneous} if $\beta_i=0$ for all $i=1,\ldots ,m,$ otherwise it is called \alert{nonhomogeneous}.

}
\frame{\frametitle{}
Collecting the coefficients in the matrix $\mathbf{A}=[\alpha_{ij}]$ the unknowns and the right-side scalars in the column vectors 
$$\mathbf{x}=\left[\begin{array}{c}
x_1\\
x_2\\
\vdots \\
x_n
\end{array}\right]\quad \mbox{and}\quad \mathbf{b}=\left[\begin{array}{c}
\beta_1\\
\beta_2\\
\vdots \\
\beta_m
\end{array}\right]$$
respectively, the matrix equation $\mathbf{Ax}=\mathbf{b}$ can be obtained.\\
Since the left side of the above matrix equation is a linear combination of the columns of $\mathbf{A}$ the following statement can be obtained:\\
\begin{theorem}
The (\ref{lsystem}) linear system of equations is solvable if and only if $\mathbf{b}$ is the linear combination of columns of the coefficient matrix $\mathbf{A}$, that is, $\mathbf{b}$ is in the subspace spanned by the columns of $\mathbf{A}$.
\end{theorem}

}


\frame{\frametitle{}
\bigskip

A solution of the system of linear equations $\mathbf{Ax}=\mathbf{b}$ is an $n$-tuple (a vector) $\mathbf{s}=(s_1,s_2,\ldots ,s_n),$ such that substituting $s_i\mbox{ for }x_i\ i=1,2,\ldots ,n$ each equation of the system (\ref{lsystem}) becomes true numerical statement. The solution set of (\ref{lsystem}) is the set of all solution vectors.
\bigskip

To provide an effective method for solution of systems of linear equation first we show a process for factoring matrices.
}
\section{Factoring matrices}
\frame{\frametitle{Basis factoring matrices}
Let $\mathbf{A}=[\mathbf{a}_1,\mathbf{a}_2,\ldots ,\mathbf{a}_n]$ be a matrix having rank $r.$ Then there is an $r$ element linearly independent set of vectors $\mathbf{a}_{i_1},\ldots,\mathbf{a}_{i_r}$ such that any vector of the subspace spanned by the columns of $\mathbf{A}$ -- which is called the {\em column space} of $\mathbf{A}$ -- can be expressed as a linear combination of the vectors $\mathbf{a}_{i_1},\ldots,\mathbf{a}_{i_r},$ that is, $\mathbf{a}_{i_1},\ldots,\mathbf{a}_{i_r}$ is a basis of the column space of $\mathbf{A}.$ Therefore
\begin{block}{}
\begin{eqnarray*}
\mathbf{a}_1&=&d_{11}\mathbf{a}_{i_1}+d_{21}\mathbf{a}_{i_2}+\cdots +d_{r1}\mathbf{a}_{i_r}\\
\mathbf{a}_2&=&d_{12}\mathbf{a}_{i_1}+d_{22}\mathbf{a}_{i_2}+\cdots +d_{r2}\mathbf{a}_{i_r}\\
&\vdots& \\
\mathbf{a}_n&=&d_{1n}\mathbf{a}_{i_1}+d_{2n}\mathbf{a}_{i_2}+\cdots +d_{rn}\mathbf{a}_{i_r}.
\end{eqnarray*}
\end{block}
}

\frame{\frametitle{Basis factoring matrices}
Form the matrices 
\begin{block}{}
\[
\mathbf{A}_1=[\mathbf{a}_{i_1},\ldots,\mathbf{a}_{i_r}] \text{ and } 
\mathbf{A}_2=\left[\begin{array}{cccc}
d_{11}&d_{12}&\ldots &d_{1n}\\
d_{21}&d_{22}&\ldots &d_{2n}\\
\vdots \\
d_{r1}&d_{r2}&\ldots &d_{rn}\\
\end{array}\right],
\]
\end{block}
then the above equation can be written in the form: 
\begin{block}{}
$$\mathbf{A}=\mathbf{A}_1\cdot \mathbf{A}_2.$$
\end{block}
Notice that the $j$th column of the matrix $\mathbf{A}_2$ is just the coordinate vector of $\mathbf{a}_j$ with respect to the basis $\{\mathbf{a}_{i_1},\ldots,\mathbf{a}_{i_r}\}.$
}

\frame{\frametitle{}
\begin{example}
Factor the matrix $\mathbf{A}=\left[\begin{array}{rrrr}
1&2&0&3\\
2&1&3&3\\
-1&1&-3&0
\end{array}\right].$
\end{example}

Using elementary basis transformation we find the rank of $\mathbf{A}$ and
a basis of its column space, starting with the standard basis
$\{\mathbf{e}_1,\mathbf{e}_2,\mathbf{e}_3\}$ of $\R^3.$
$$\begin{array}{c|rrrr}
&\mathbf{a}_1&\mathbf{a}_2&\mathbf{a}_3&\mathbf{a}_4\\ \hline
\mathbf{e}_1&\fbox{$1$}&2&0&3\\
\mathbf{e}_2&2&1&3&3\\
\mathbf{e}_3&-1&1&-3&0
\end{array}\longrightarrow \begin{array}{c|rrr}
&\mathbf{a}_2&\mathbf{a}_3&\mathbf{a}_4\\ \hline
\mathbf{a}_1&2&0&3\\
\mathbf{e}_2&\fbox{$-3$}&3&-3\\
\mathbf{e}_3&3&-3&3
\end{array}\longrightarrow \begin{array}{c|rr}
&\mathbf{a}_3&\mathbf{a}_4\\ \hline
\mathbf{a}_1&2&1\\
\mathbf{a}_2&-1&1\\
\mathbf{e}_3&0&0
\end{array}$$
It can be red off from the last table, that the first and second column of $\mathbf{A}$ form a basis of its column space and $\mathbf{a}_3=2\mathbf{a}_1-\mathbf{a}_2$ and $\mathbf{a}_4=\mathbf{a}_1+\mathbf{a}_2.$ 
}
\frame{\frametitle{}
Remember that $j$-th column of the matrix $\mathbf{A_2}$ is the coordinate vector of $\mathbf{a_j}$ with respect to the 
basis $\left\{ \mathbf{a_1},\mathbf{a_2} \right\}$.
Therefore
\begin{exampleblock}{}
$$\mathbf{A}_1=[\mathbf{a}_1,\mathbf{a}_2]=\left[\begin{array}{rr}
1&2\\
2&1\\
-1&1
\end{array}\right]\quad\text{ and }\quad \mathbf{A}_2=\left[\begin{array}{rrrr}
1&0&2&1\\
0&1&-1&1
\end{array}\right],$$
\end{exampleblock}
where we used the obvious fact that $\left[\begin{array}{c}
1\\
0
\end{array}\right]$ and $\left[\begin{array}{c}
0\\
1
\end{array}\right]$ are the coordinate vectors of $\mathbf{a}_1$ respectively $\mathbf{a}_2$ with respect to the basis $\{\mathbf{a}_1,\mathbf{a}_2\}.$

}

\section{The rank theorem}
\begin{frame}
Let $A$ be an $m\times n$ matrix, and consider the factorization $A=B\cdot C$ as before.
Here $B$ includes a maximal linearly independent system of columns of $A$. 
Because of the maximal property the span of columns of $B$ coincides with the column space of $A$.
It means that the columns of $B$ form a basis of the column space of $A$.

We know that $C$ is the last elimination table after dropping the zero rows.
The system of rows of $C$ is linearly independent, because every row has pivot position $1$,
and at that position the rest of the rows includes $0$. 
Thus none of the rows can be obtained as the linear combination of the rest.
One can see that the row space is not changing during the elimination.
It means that the rows of $C$ form a basis of the column space of $A$.

Thus we proved the 
    \begin{theorem}[Rank theorem]
        The dimensions of the column space and the row space of a matrix are equal.
        This common dimension is called the \alert{rank of the matrix}.
        The rank is the same as the number of the pivot terms of the Gauss-Jordan elimination of the matrix.
    \end{theorem}
\end{frame}
\section{Problems}
\begin{frame}
    \frametitle{Problems}

    \begin{exampleblock}{Factorize the matrices below}
    \begin{columns}[t]
    \column{.5\textwidth}
    \begin{itemize}[<+|alert@+>]
        \item
            \(
            \left[
            \begin{array}{rrr}
                1&3&5\\
                2&4&6\\
                3&5&7
            \end{array}
            \right]
            \)
        \item
            \(
            \left[
            \begin{array}{rrr}
                1&-3&2\\
                2&-5&5\\
                1&-2&3
            \end{array}
            \right]
            \)
        \item
            \(
            \left[
            \begin{array}{rrr}
                -3&-5&36\\
                -1&0&7\\
                1&1&-10
            \end{array}
            \right]
            \)
    \end{itemize}
    \column{.5\textwidth}
    \begin{itemize}[<+|alert@+>]
        \item
            \(
            \left[
            \begin{array}{rrrr}
                1&2&-3&-4\\
                1&3&-3&-4\\
                2&2&-6&-8
            \end{array}
            \right]
            \)
        \item
            \(
            \left[
            \begin{array}{rrrr}
                4&8&-4&4\\
                3&8&5&-11\\
                -2&1&12&-17
            \end{array}
            \right]
            \)
    \end{itemize}
    \end{columns}
    \end{exampleblock}
\end{frame}
\end{document}
Created:        Sun 27 Apr 2008 11:32:30 PM CEST
Last Modified:  Tue 30 Mar 2010 07:30:19 AM CEST
\frame{\frametitle{Polytope}
The {\em polytope} is the convex hull of a finite set of vectors can be considered as the generalization of convex polygons and solids like triangle, rectangle, tetrahedron and so forth.

The generalization of the plane is given as follows:
\begin{definition}
Let $\mathbf{n}\in \R^n$ and $\alpha\in \R$ such that $\mathbf{n}\neq \mathbf{0}.$ Then the hyperplane $S(\mathbf{n},\alpha )$ is the set of all vectors $\mathbf{x}\in \R^n$ satisfying the equation $\left\langle \mathbf{n};\mathbf{x}\right\rangle =\alpha .$
\end{definition}

For any two points $\mathbf{x}_1,\mathbf{x}_2\in S(\mathbf{n},\alpha )$ the vector $\mathbf{x}_1-\mathbf{x}_2$ is orthogonal to the vector $\mathbf{n},$ therefore it is often called the normal vector of the hyperplane.

}
